\section*{习题 6.6\quad 定积分的应用}

1. 求下列曲线的弧长

1. 题目：\(y=x^{\frac{3}{2}},\ x\in[0,4]\)

解答：弧长公式：\(L=\int_{a}^{b}\sqrt{1+(y')^{2}}\,dx\)。

计算导数：\(y'=\frac{3}{2}x^{\frac{1}{2}}\)。
\[ L=\int_{0}^{4}\sqrt{1+\left(\frac{3}{2}\sqrt{x}\right)^{2}}dx=\int_{0}^{4}\sqrt{1+\frac{9}{4}x}\,dx. \]
令 \(u=1+\frac{9}{4}x\)，则 \(du=\frac{9}{4}dx,\ dx=\frac{4}{9}du\)。当 \(x=0\) 时 \(u=1\)；当 \(x=4\) 时 \(u=1+\frac{9}{4}\times4=10\)。
\[ L=\int_{1}^{10}\sqrt{u}\cdot\frac{4}{9}du=\frac{4}{9}\left[\frac{2}{3}u^{\frac{3}{2}}\right]_{1}^{10}=\frac{8}{27}\left(10^{\frac{3}{2}}-1\right)=\frac{8}{27}\left(10\sqrt{10}-1\right). \]
答案：\(\frac{8}{27}\left(10\sqrt{10}-1\right)\)

2. 题目：\(x=\frac{1}{4}y^{2}-\frac{1}{2}\ln y,\ y\in[1,e]\)

解答：弧长公式：\(L=\int_{a}^{b}\sqrt{1+(x')^{2}}\,dy\)。

计算导数：\(x'=\frac{1}{2}y-\frac{1}{2y}=\frac{y^{2}-1}{2y}\)。
\[ L=\int_{1}^{e}\sqrt{1+\left(\frac{y^{2}-1}{2y}\right)^{2}}dy=\int_{1}^{e}\sqrt{1+\frac{(y^{2}-1)^{2}}{4y^{2}}}\,dy=\int_{1}^{e}\sqrt{\frac{4y^{2}+(y^{2}-1)^{2}}{4y^{2}}}\,dy \]
\[ =\int_{1}^{e}\frac{\sqrt{4y^{2}+y^{4}-2y^{2}+1}}{2y}dy=\int_{1}^{e}\frac{\sqrt{y^{4}+2y^{2}+1}}{2y}dy=\int_{1}^{e}\frac{\sqrt{(y^{2}+1)^{2}}}{2y}dy=\int_{1}^{e}\frac{y^{2}+1}{2y}dy \]
\[ =\frac{1}{2}\int_{1}^{e}\left(y+\frac{1}{y}\right)dy=\frac{1}{2}\left[\frac{y^{2}}{2}+\ln y\right]_{1}^{e}=\frac{1}{2}\left(\frac{e^{2}}{2}+1-\frac{1}{2}-0\right)=\frac{e^{2}+1}{4}. \]
答案：\(\frac{e^{2}+1}{4}\)

3. 题目：\(\sqrt{x}+\sqrt{y}=1\)

解答：由方程得 \(y=(1-\sqrt{x})^{2}=1-2\sqrt{x}+x,\ x\in[0,1]\)。

计算导数：\(y'=1-\frac{1}{\sqrt{x}}\)。
\[ L=\int_{0}^{1}\sqrt{1+\left(1-\frac{1}{\sqrt{x}}\right)^{2}}dx=\int_{0}^{1}\sqrt{1+1-\frac{2}{\sqrt{x}}+\frac{1}{x}}\,dx=\int_{0}^{1}\sqrt{\frac{1}{x}-\frac{2}{\sqrt{x}}+2}\,dx=\int_{0}^{1}\sqrt{\left(\frac{1}{\sqrt{x}}-1\right)^{2}+1}\,dx. \]
令 \(u=\sqrt{x}\)，则 \(x=u^{2},\ dx=2u\,du\)；当 \(x=0\) 时 \(u=0\)，当 \(x=1\) 时 \(u=1\)。
\[ L=\int_{0}^{1}\sqrt{\left(\frac{1}{u}-1\right)^{2}+1}\cdot2u\,du=2\int_{0}^{1}\sqrt{\frac{(1-u)^{2}}{u^{2}}+1}\cdot u\,du=2\int_{0}^{1}\sqrt{(1-u)^{2}+u^{2}}\,du \]
\[ =2\int_{0}^{1}\sqrt{1-2u+u^{2}+u^{2}}\,du=2\int_{0}^{1}\sqrt{2u^{2}-2u+1}\,du. \]
完成平方：
\[ 2u^{2}-2u+1=2(u^{2}-u)+1=2\left[\left(u-\frac{1}{2}\right)^{2}-\frac{1}{4}\right]+1=2\left(u-\frac{1}{2}\right)^{2}+\frac{1}{2}. \]
令 \(t=\sqrt{2}\left(u-\frac{1}{2}\right)\)，则 \(du=\frac{dt}{\sqrt2}\)；当 \(u=0\) 时 \(t=-\frac{\sqrt2}{2}\)，当 \(u=1\) 时 \(t=\frac{\sqrt2}{2}\)。
\[ L=2\int_{-\frac{\sqrt2}{2}}^{\frac{\sqrt2}{2}}\sqrt{t^{2}+\frac{1}{2}}\cdot\frac{dt}{\sqrt2}=\sqrt2\int_{-\frac{\sqrt2}{2}}^{\frac{\sqrt2}{2}}\sqrt{t^{2}+\frac{1}{2}}\,dt. \]
由于被积函数是偶函数：
\[ L=2\sqrt2\int_{0}^{\frac{\sqrt2}{2}}\sqrt{t^{2}+\frac{1}{2}}\,dt. \]
使用公式 \(\int\sqrt{t^{2}+a^{2}}\,dt=\frac{t}{2}\sqrt{t^{2}+a^{2}}+\frac{a^{2}}{2}\ln\left|t+\sqrt{t^{2}+a^{2}}\right|+C\)，其中 \(a^{2}=\frac{1}{2}\)：
\[ L=2\sqrt2\left[\frac{t}{2}\sqrt{t^{2}+\frac{1}{2}}+\frac{1}{4}\ln\left|t+\sqrt{t^{2}+\frac{1}{2}}\right|\right]_{0}^{\frac{\sqrt2}{2}}. \]
当 \(t=\frac{\sqrt2}{2}\) 时：\(t^{2}+\frac{1}{2}=\frac{1}{2}+\frac{1}{2}=1\)，\(\sqrt{t^{2}+\frac{1}{2}}=1\)，
\[ \frac{t}{2}\sqrt{t^{2}+\frac{1}{2}}=\frac{\sqrt2}{4}\cdot1=\frac{\sqrt2}{4},\qquad \ln\left|t+\sqrt{t^{2}+\frac{1}{2}}\right|=\ln\left(1+\frac{\sqrt2}{2}\right). \]
当 \(t=0\) 时：
\[ \frac{t}{2}\sqrt{t^{2}+\frac{1}{2}}=0,\qquad \ln\left|0+\sqrt{0+\frac{1}{2}}\right|=\ln\frac{\sqrt2}{2}=-\frac{1}{2}\ln2. \]
所以
\[ L=2\sqrt2\left[\frac{\sqrt2}{4}+\frac{1}{4}\ln\left(1+\frac{\sqrt2}{2}\right)-\left(0+\frac{1}{4}\left(-\frac{1}{2}\ln2\right)\right)\right] \]
\[ =2\sqrt2\left[\frac{\sqrt2}{4}+\frac{1}{4}\ln\left(1+\frac{\sqrt2}{2}\right)+\frac{1}{8}\ln2\right]=1+\frac{\sqrt2}{2}\ln\left(1+\frac{\sqrt2}{2}\right)+\frac{\sqrt2}{4}\ln2. \]
化简对数项：\(\ln\left(1+\frac{\sqrt2}{2}\right)=\ln\left(\frac{2+\sqrt2}{2}\right)=\ln(2+\sqrt2)-\ln2\)，于是
\[ \frac{\sqrt2}{2}\ln\left(1+\frac{\sqrt2}{2}\right)+\frac{\sqrt2}{4}\ln2=\frac{\sqrt2}{2}\left[\ln(2+\sqrt2)-\ln2\right]+\frac{\sqrt2}{4}\ln2=\frac{\sqrt2}{2}\ln(2+\sqrt2)-\frac{\sqrt2}{4}\ln2. \]
所以 \(L=1+\frac{\sqrt2}{2}\ln(2+\sqrt2)-\frac{\sqrt2}{4}\ln2\)。

进一步简化：\(2+\sqrt2=\sqrt2(1+\sqrt2)\)，所以 \(\ln(2+\sqrt2)=\ln\sqrt2+\ln(1+\sqrt2)=\frac{1}{2}\ln2+\ln(1+\sqrt2)\)。代入：
\[ L=1+\frac{\sqrt2}{2}\left(\frac{1}{2}\ln2+\ln(1+\sqrt2)\right)-\frac{\sqrt2}{4}\ln2=1+\frac{\sqrt2}{4}\ln2+\frac{\sqrt2}{2}\ln(1+\sqrt2)-\frac{\sqrt2}{4}\ln2=1+\frac{\sqrt2}{2}\ln(1+\sqrt2). \]
答案：\(1+\frac{\sqrt2}{2}\ln(1+\sqrt2)\)

4. 题目：\(x^{\frac{2}{3}}+y^{\frac{2}{3}}=a^{\frac{2}{3}}\)

解答：这是星形线，使用参数方程：\(x=a\cos^{3}t,\ y=a\sin^{3}t,\ t\in[0,2\pi]\)。

弧长公式：\(L=\int_{0}^{2\pi}\sqrt{(x')^{2}+(y')^{2}}\,dt\)，其中 \(x'=-3a\cos^{2}t\sin t,\ y'=3a\sin^{2}t\cos t\)。
\[ L=\int_{0}^{2\pi}\sqrt{9a^{2}\cos^{4}t\sin^{2}t+9a^{2}\sin^{4}t\cos^{2}t}\,dt=\int_{0}^{2\pi}3a|\sin t\cos t|\sqrt{\cos^{2}t+\sin^{2}t}\,dt \]
\[ =3a\int_{0}^{2\pi}|\sin t\cos t|\,dt=3a\cdot4\int_{0}^{\frac{\pi}{2}}\sin t\cos t\,dt=12a\left[\frac{\sin^{2}t}{2}\right]_{0}^{\frac{\pi}{2}}=12a\cdot\frac{1}{2}=6a. \]
答案：\(6a\)

5. 题目：\(\begin{cases}x=a(\cos t+t\sin t),\\ y=a(\sin t-t\cos t),\end{cases}\ t\in[0,2\pi]\)

解答：计算导数：
\[ x'=a(-\sin t+\sin t+t\cos t)=at\cos t,\qquad y'=a(\cos t-\cos t+t\sin t)=at\sin t. \]
\[ L=\int_{0}^{2\pi}\sqrt{(at\cos t)^{2}+(at\sin t)^{2}}\,dt=\int_{0}^{2\pi}at\,dt=a\left[\frac{t^{2}}{2}\right]_{0}^{2\pi}=2a\pi^{2}. \]
答案：\(2a\pi^{2}\)

6. 题目：\(r=a(1+\cos\theta)\)（心形线）

解答：极坐标下弧长公式：\(L=\int_{\theta_{1}}^{\theta_{2}}\sqrt{r^{2}+(r')^{2}}\,d\theta\)，其中 \(r'=-a\sin\theta\)。
\[ L=\int_{0}^{2\pi}\sqrt{a^{2}(1+\cos\theta)^{2}+a^{2}\sin^{2}\theta}\,d\theta=a\int_{0}^{2\pi}\sqrt{1+2\cos\theta+\cos^{2}\theta+\sin^{2}\theta}\,d\theta \]
\[ =a\int_{0}^{2\pi}\sqrt{2+2\cos\theta}\,d\theta=a\int_{0}^{2\pi}\sqrt{4\cos^{2}\frac{\theta}{2}}\,d\theta=2a\int_{0}^{2\pi}\left|\cos\frac{\theta}{2}\right|d\theta. \]
在 \([0,2\pi]\) 上，\(\cos\frac{\theta}{2}\) 在 \([0,\pi]\) 非负，在 \([\pi,2\pi]\) 非正，因此
\[ L=2a\left(\int_{0}^{\pi}\cos\frac{\theta}{2}\,d\theta+\int_{\pi}^{2\pi}\left(-\cos\frac{\theta}{2}\right)d\theta\right)=2a\left(\left[2\sin\frac{\theta}{2}\right]_{0}^{\pi}-\left[2\sin\frac{\theta}{2}\right]_{\pi}^{2\pi}\right)=4a(1-0-0+1)=8a. \]
答案：\(8a\)

2. 填空题

1. 题目：设 \(f(x),g(x)\) 在 \([a,b]\) 上连续，曲线 \(y=f(x),\ y=g(x)\) 及直线 \(x=a,\ x=b\) 所围成平面区域的面积为 \_\_\_。

解答：面积公式：\(S=\int_{a}^{b}|f(x)-g(x)|\,dx\)。

答案：\(\int_{a}^{b}|f(x)-g(x)|\,dx\)

2. 题目：设 \(f(x),g(x)\) 在 \([a,b]\) 上连续，曲线 \(y=f(x),\ y=g(x)\) 及直线 \(x=a,\ x=b\) 所围成的平面区域绕 \(x\) 轴旋转一周所得旋转体的体积为 \_\_\_。

解答：旋转体体积公式（圆盘法）：\(V=\pi\int_{a}^{b}\left|[f(x)]^{2}-[g(x)]^{2}\right|dx\)。

答案：\(\pi\int_{a}^{b}\left|[f(x)]^{2}-[g(x)]^{2}\right|dx\)

3. 在直角坐标系下，求下列曲线所围成平面图形的面积

1. 题目：\(y=x^{2},\ x+y=2\)

解答：求交点：\(x^{2}+x-2=0\Rightarrow(x+2)(x-1)=0\Rightarrow x=-2,1\)。

面积：
\[ S=\int_{-2}^{1}\left[(2-x)-x^{2}\right]dx=\int_{-2}^{1}(2-x-x^{2})dx=\left[2x-\frac{x^{2}}{2}-\frac{x^{3}}{3}\right]_{-2}^{1} \]
\[ =\left(2-\frac{1}{2}-\frac{1}{3}\right)-\left(-4-2+\frac{8}{3}\right)=\frac{7}{6}-\left(-\frac{10}{3}\right)=\frac{9}{2}. \]
答案：\(\frac{9}{2}\)

2. 题目：\(y=\ln x,\ y=0\) 及 \(x=\frac{1}{e},\ x=e\)

解答：面积：\(S=\int_{\frac{1}{e}}^{e}|\ln x|\,dx\)。在 \(\left[\frac{1}{e},1\right]\) 上 \(\ln x\le0\)，在 \([1,e]\) 上 \(\ln x\ge0\)，所以
\[ S=\int_{\frac{1}{e}}^{1}(-\ln x)dx+\int_{1}^{e}\ln x\,dx=\left[-x\ln x+x\right]_{\frac{1}{e}}^{1}+\left[x\ln x-x\right]_{1}^{e} \]
\[ =\left(0+1-\left(\frac{1}{e}+\frac{1}{e}\right)\right)+(e-e-(0-1))=2-\frac{2}{e}. \]
答案：\(2-\frac{2}{e}\)

3. 题目：\(\frac{x^{2}}{a^{2}}+\frac{y^{2}}{b^{2}}=1\)

解答：椭圆面积：\(S=\pi ab\)。

答案：\(\pi ab\)

4. 题目：\(y=e^{x},\ y=e^{-x},\ x=1\)

解答：求交点：\(e^{x}=e^{-x}\Rightarrow x=0\)。

面积：
\[ S=\int_{0}^{1}(e^{x}-e^{-x})dx=\left[e^{x}+e^{-x}\right]_{0}^{1}=(e+e^{-1})-(1+1)=e+\frac{1}{e}-2. \]
答案：\(e+\frac{1}{e}-2\)

5. 题目：\(y=\sin x,\ y=\cos x,\ x=0,\ x=\frac{\pi}{2}\)

解答：求交点：\(\sin x=\cos x\Rightarrow x=\frac{\pi}{4}\)。

面积：
\[ S=\int_{0}^{\frac{\pi}{4}}(\cos x-\sin x)dx+\int_{\frac{\pi}{4}}^{\frac{\pi}{2}}(\sin x-\cos x)dx \]
\[ =\left[\sin x+\cos x\right]_{0}^{\frac{\pi}{4}}+\left[-\cos x-\sin x\right]_{\frac{\pi}{4}}^{\frac{\pi}{2}}=\left(\sqrt2-1\right)+\left(-1+\sqrt2\right)=2\sqrt2-2. \]
答案：\(2\sqrt2-2\)

6. 题目：圆 \(x^{2}+y^{2}=8\) 被抛物线 \(y=\frac{1}{2}x^{2}\) 分成的两部分

解答：求交点：\(x^{2}+\left(\frac{1}{2}x^{2}\right)^{2}=8\Rightarrow x^{2}+\frac{1}{4}x^{4}=8\Rightarrow x^{4}+4x^{2}-32=0\)。

令 \(t=x^{2}\)，则 \(t^{2}+4t-32=0\Rightarrow t=4\)（负根舍去）\(\Rightarrow x=\pm2\)。

由对称性，计算右半部分面积。圆在抛物线之上的部分：
\[ S_{1}=\int_{0}^{2}\left(\sqrt{8-x^{2}}-\frac{1}{2}x^{2}\right)dx. \]
圆在抛物线之下的部分：\(S_{2}=\frac{1}{2}\pi(\sqrt8)^{2}-S_{1}=4\pi-S_{1}\)。

计算 \(S_{1}\)：
\[ \int_{0}^{2}\sqrt{8-x^{2}}dx=\int_{0}^{\frac{\pi}{4}}\sqrt{8-8\sin^{2}\theta}\cdot2\sqrt2\cos\theta\,d\theta=\int_{0}^{\frac{\pi}{4}}8\cos^{2}\theta\,d\theta=4\int_{0}^{\frac{\pi}{4}}(1+\cos2\theta)d\theta \]
\[ =4\left[\theta+\frac{\sin2\theta}{2}\right]_{0}^{\frac{\pi}{4}}=4\left(\frac{\pi}{4}+\frac{1}{2}\right)=\pi+2, \]
\[ \int_{0}^{2}\frac{1}{2}x^{2}dx=\frac{1}{2}\left[\frac{x^{3}}{3}\right]_{0}^{2}=\frac{8}{6}=\frac{4}{3}. \]
所以 \(S_{1}=\pi+2-\frac{4}{3}=\pi+\frac{2}{3}\)，从而 \(S_{2}=4\pi-\left(\pi+\frac{2}{3}\right)=3\pi-\frac{2}{3}\)。

整个圆被分成的两部分面积（对称性）：在抛物线之上的部分为 \(2S_{1}=2\pi+\frac{4}{3}\)，在抛物线之下的部分为 \(2S_{2}=6\pi-\frac{4}{3}\)。

答案：两部分面积分别为 \(2\pi+\frac{4}{3}\) 和 \(6\pi-\frac{4}{3}\)

4. 求由极坐标方程或参数方程表示的曲线所围成的平面图形的面积

1. 题目：\(r=2\cos\theta\)

解答：这是圆，极坐标下面积公式：\(S=\frac{1}{2}\int_{\theta_{1}}^{\theta_{2}}r^{2}d\theta\)。\(r=2\cos\theta\) 在 \(\theta\in\left[-\frac{\pi}{2},\frac{\pi}{2}\right]\) 时 \(r\ge0\)。
\[ S=\frac{1}{2}\int_{-\frac{\pi}{2}}^{\frac{\pi}{2}}(2\cos\theta)^{2}d\theta=2\int_{-\frac{\pi}{2}}^{\frac{\pi}{2}}\cos^{2}\theta\,d\theta=2\int_{-\frac{\pi}{2}}^{\frac{\pi}{2}}\frac{1+\cos2\theta}{2}d\theta=\int_{-\frac{\pi}{2}}^{\frac{\pi}{2}}(1+\cos2\theta)d\theta \]
\[ =\left[\theta+\frac{\sin2\theta}{2}\right]_{-\frac{\pi}{2}}^{\frac{\pi}{2}}=\left(\frac{\pi}{2}+0\right)-\left(-\frac{\pi}{2}+0\right)=\pi. \]
答案：\(\pi\)

2. 题目：\(r^{2}=2a^{2}\cos2\theta\)

解答：这是双纽线，面积：
\[ S=4\times\frac{1}{2}\int_{0}^{\frac{\pi}{4}}2a^{2}\cos2\theta\,d\theta=4a^{2}\int_{0}^{\frac{\pi}{4}}\cos2\theta\,d\theta=4a^{2}\left[\frac{\sin2\theta}{2}\right]_{0}^{\frac{\pi}{4}}=2a^{2}. \]
答案：\(2a^{2}\)

3. 题目：\(x=a\cos^{3}t,\ y=a\sin^{3}t\)（星形线）

解答：参数方程下面积公式：
\[ S=\left|\oint y\,dx\right|=\left|\int y(t)x'(t)dt\right|, \]
其中 \(x'=-3a\cos^{2}t\sin t\)。
\[ S=\left|\int_{0}^{2\pi}a\sin^{3}t\cdot(-3a\cos^{2}t\sin t)dt\right|=3a^{2}\int_{0}^{2\pi}\sin^{4}t\cos^{2}t\,dt. \]
由对称性：
\[ S=4\times3a^{2}\int_{0}^{\frac{\pi}{2}}\sin^{4}t\cos^{2}t\,dt=12a^{2}\int_{0}^{\frac{\pi}{2}}\sin^{4}t(1-\sin^{2}t)dt=12a^{2}\left(\int_{0}^{\frac{\pi}{2}}\sin^{4}t\,dt-\int_{0}^{\frac{\pi}{2}}\sin^{6}t\,dt\right). \]
使用华里士公式：
\[ \int_{0}^{\frac{\pi}{2}}\sin^{n}t\,dt=\frac{(n-1)!!}{n!!}\cdot\frac{\pi}{2},\qquad \int_{0}^{\frac{\pi}{2}}\sin^{n}t\,dt=\frac{(n-1)!!}{n!!}, \]
其中第一式对偶数 \(n\) 成立，第二式对奇数 \(n\) 成立。
\[ \int_{0}^{\frac{\pi}{2}}\sin^{4}t\,dt=\frac{3!!}{4!!}\cdot\frac{\pi}{2}=\frac{3}{4}\cdot\frac{1}{2}\cdot\frac{\pi}{2}=\frac{3\pi}{16}, \]
\[ \int_{0}^{\frac{\pi}{2}}\sin^{6}t\,dt=\frac{5!!}{6!!}\cdot\frac{\pi}{2}=\frac{5\cdot3\cdot1}{6\cdot4\cdot2}\cdot\frac{\pi}{2}=\frac{15}{48}\cdot\frac{\pi}{2}=\frac{5\pi}{32}. \]
所以
\[ S=12a^{2}\left(\frac{3\pi}{16}-\frac{5\pi}{32}\right)=12a^{2}\cdot\frac{\pi}{32}=\frac{3\pi a^{2}}{8}. \]
答案：\(\frac{3\pi a^{2}}{8}\)

5. 求下列已知曲线所围的图形，绕指定的坐标轴旋转而形成的旋转体的体积

1. 题目：\(y=x^{2},\ y=\sqrt{x}\)，绕 \(y\) 轴

解答：求交点：\(x^{2}=\sqrt{x}\Rightarrow x^{4}=x\Rightarrow x(x^{3}-1)=0\Rightarrow x=0,1\)。

使用柱壳法：
\[ V=2\pi\int_{0}^{1}x|\sqrt{x}-x^{2}|\,dx=2\pi\int_{0}^{1}(x^{\frac{3}{2}}-x^{3})dx=2\pi\left[\frac{2}{5}x^{\frac{5}{2}}-\frac{x^{4}}{4}\right]_{0}^{1}=2\pi\left(\frac{2}{5}-\frac{1}{4}\right)=2\pi\cdot\frac{3}{20}=\frac{3\pi}{10}. \]
答案：\(\frac{3\pi}{10}\)

2. 题目：\(x^{2}+(y-b)^{2}=a^{2}\ (0<a\le b)\)，绕 \(x\) 轴

解答：这是圆心在 \((0,b)\)、半径为 \(a\) 的圆。使用圆盘法：
\[ V=\pi\int_{-a}^{a}\left[\left(b+\sqrt{a^{2}-x^{2}}\right)^{2}-\left(b-\sqrt{a^{2}-x^{2}}\right)^{2}\right]dx=\pi\int_{-a}^{a}4b\sqrt{a^{2}-x^{2}}\,dx=8\pi b\int_{0}^{a}\sqrt{a^{2}-x^{2}}\,dx. \]
令 \(x=a\sin t\)，则 \(dx=a\cos t\,dt\)，
\[ V=8\pi b\int_{0}^{\frac{\pi}{2}}a\cos t\cdot a\cos t\,dt=8\pi a^{2}b\int_{0}^{\frac{\pi}{2}}\cos^{2}t\,dt=8\pi a^{2}b\cdot\frac{\pi}{4}=2\pi^{2}a^{2}b. \]
答案：\(2\pi^{2}a^{2}b\)

3. 题目：\(\begin{cases}x=a\cos^{3}t,\\ y=a\sin^{3}t,\end{cases}\) 绕 \(x\) 轴

解答：星形线绕 \(x\) 轴旋转的体积。使用圆盘法：
\[ V=\pi\int_{-a}^{a}y^{2}dx=\pi\int_{\pi}^{0}(a\sin^{3}t)^{2}\cdot(-3a\cos^{2}t\sin t)dt=3\pi a^{3}\int_{0}^{\pi}\sin^{7}t\cos^{2}t\,dt=6\pi a^{3}\int_{0}^{\frac{\pi}{2}}\sin^{7}t(1-\sin^{2}t)dt \]
\[ =6\pi a^{3}\left(\int_{0}^{\frac{\pi}{2}}\sin^{7}t\,dt-\int_{0}^{\frac{\pi}{2}}\sin^{9}t\,dt\right). \]
使用公式 \(\int_{0}^{\frac{\pi}{2}}\sin^{n}t\,dt=\frac{(n-1)!!}{n!!}\cdot\frac{\pi}{2}\)（\(n\) 为偶数）或 \(\frac{(n-1)!!}{n!!}\)（\(n\) 为奇数）：
\[ \int_{0}^{\frac{\pi}{2}}\sin^{7}t\,dt=\frac{6!!}{7!!}=\frac{6\cdot4\cdot2}{7\cdot5\cdot3\cdot1}=\frac{48}{105}=\frac{16}{35}, \]
\[ \int_{0}^{\frac{\pi}{2}}\sin^{9}t\,dt=\frac{8!!}{9!!}=\frac{8\cdot6\cdot4\cdot2}{9\cdot7\cdot5\cdot3\cdot1}=\frac{384}{945}=\frac{128}{315}. \]
所以
\[ V=6\pi a^{3}\left(\frac{16}{35}-\frac{128}{315}\right)=6\pi a^{3}\cdot\frac{16}{315}=\frac{32\pi a^{3}}{105}. \]
答案：\(\frac{32\pi a^{3}}{105}\)

6. 旋转曲面的侧面积

证明题：设 \(f(x)\) 在 \([a,b]\) 上连续可导，且 \(f(x)\ge0\)。证明曲线段 \(y=f(x),\ x\in[a,b]\) 绕 \(x\) 轴旋转一周所得旋转曲面的面积为
\[ S=2\pi\int_{a}^{b}f(x)\sqrt{1+(f'(x))^{2}}\,dx. \]

证明：在 \([x,x+dx]\) 小区间上，曲线段近似为直线段，长度 \(ds=\sqrt{1+(f'(x))^{2}}\,dx\)，旋转形成的曲面近似为圆台侧面积。以 \(f(x)\) 为平均半径、以 \(ds\) 为斜高，得
\[ dS=2\pi f(x)ds=2\pi f(x)\sqrt{1+(f'(x))^{2}}\,dx. \]
积分得
\[ S=2\pi\int_{a}^{b}f(x)\sqrt{1+(f'(x))^{2}}\,dx. \]

1. 题目：\(y=2\sqrt{x},\ x\in[0,1]\)，绕 \(x\) 轴

解答：\(y'=\frac{1}{\sqrt{x}}\)。
\[ S=2\pi\int_{0}^{1}2\sqrt{x}\sqrt{1+\frac{1}{x}}\,dx=4\pi\int_{0}^{1}\sqrt{x+1}\,dx=4\pi\left[\frac{2}{3}(x+1)^{\frac{3}{2}}\right]_{0}^{1}=\frac{8\pi}{3}\left(2\sqrt2-1\right). \]
答案：\(\frac{8\pi}{3}\left(2\sqrt2-1\right)\)

2. 题目：\(x^{\frac{3}{2}}+y^{\frac{3}{2}}=a^{\frac{3}{2}}\)，绕 \(x\) 轴

解答：由方程得 \(y=\left(a^{\frac{3}{2}}-x^{\frac{3}{2}}\right)^{\frac{2}{3}}\)。
\[ y'=\frac{2}{3}\left(a^{\frac{3}{2}}-x^{\frac{3}{2}}\right)^{-\frac{1}{3}}\cdot\left(-\frac{3}{2}x^{\frac{1}{2}}\right)=-\frac{\sqrt{x}}{\sqrt[3]{a^{\frac{3}{2}}-x^{\frac{3}{2}}}}. \]
旋转曲面面积：
\[ S=2\pi\int_{0}^{a}y\sqrt{1+(y')^{2}}\,dx. \]
令 \(u=\left(\frac{x}{a}\right)^{\frac{3}{2}}\)，则 \(x=au^{\frac{2}{3}},\ dx=a\cdot\frac{2}{3}u^{-\frac{1}{3}}du\)。当 \(x=0\) 时 \(u=0\)，当 \(x=a\) 时 \(u=1\)。此时
\[ y=\left(a^{\frac{3}{2}}-a^{\frac{3}{2}}u\right)^{\frac{2}{3}}=a(1-u)^{\frac{2}{3}}, \]
\[ y'=-\frac{\sqrt{a}u^{\frac{1}{3}}}{a^{\frac{1}{2}}(1-u)^{\frac{1}{3}}}=-\frac{u^{\frac{1}{3}}}{(1-u)^{\frac{1}{3}}}. \]
于是
\[ 1+(y')^{2}=1+\frac{u^{\frac{2}{3}}}{(1-u)^{\frac{2}{3}}}=\frac{(1-u)^{\frac{2}{3}}+u^{\frac{2}{3}}}{(1-u)^{\frac{2}{3}}},\qquad \sqrt{1+(y')^{2}}=\frac{\sqrt{(1-u)^{\frac{2}{3}}+u^{\frac{2}{3}}}}{(1-u)^{\frac{1}{3}}}. \]
\[ S=2\pi\int_{0}^{1}a(1-u)^{\frac{2}{3}}\cdot\frac{\sqrt{(1-u)^{\frac{2}{3}}+u^{\frac{2}{3}}}}{(1-u)^{\frac{1}{3}}}\cdot a\cdot\frac{2}{3}u^{-\frac{1}{3}}du=\frac{4\pi a^{2}}{3}\int_{0}^{1}(1-u)^{\frac{1}{3}}\sqrt{(1-u)^{\frac{2}{3}}+u^{\frac{2}{3}}}\,u^{-\frac{1}{3}}du. \]
这是一个复杂的积分，数值计算可得：
\[ S\approx5.505393269817487a^{2}. \]
答案：\(S\approx5.505393269817487a^{2}\)

7. 等腰梯形薄板所受的压力

题目：高为 \(h\) 而底分别为 \(a\) 和 \(b\) 的等腰梯形板垂直悬在密度为 \(\rho\) 的液体中，薄板顶端到液面的距离为 \(c\)。求薄板所受的压力。

解答：建立坐标系：以液面为 \(y=0\)，向下为 \(y\) 轴正方向。梯形上底 \(a\)（在 \(y=c\) 处），下底 \(b\)（在 \(y=c+h\) 处）。

在深度 \(y\) 处，梯形宽度为
\[ w(y)=a+\frac{b-a}{h}(y-c). \]
压力
\[ P=\rho g\int_{c}^{c+h}y\cdot w(y)\,dy=\rho g\int_{c}^{c+h}y\left[a+\frac{b-a}{h}(y-c)\right]dy. \]
令 \(t=y-c\)，则 \(t\in[0,h]\)。
\[ P=\rho g\int_{0}^{h}(t+c)\left[a+\frac{b-a}{h}t\right]dt=\rho g\int_{0}^{h}\left[a(t+c)+\frac{b-a}{h}t(t+c)\right]dt \]
\[ =\rho g\left[a\left(\frac{t^{2}}{2}+ct\right)+\frac{b-a}{h}\left(\frac{t^{3}}{3}+\frac{ct^{2}}{2}\right)\right]_{0}^{h}=\rho g\left[a\left(\frac{h^{2}}{2}+ch\right)+\frac{b-a}{h}\left(\frac{h^{3}}{3}+\frac{ch^{2}}{2}\right)\right] \]
\[ =\rho g\left[\frac{ah^{2}}{2}+ach+\frac{(b-a)h^{2}}{3}+\frac{(b-a)ch}{2}\right]=\rho gh\left[\frac{ah}{2}+ac+\frac{(b-a)h}{3}+\frac{(b-a)c}{2}\right] \]
\[ =\rho gh\left[ac+\frac{ah}{2}+\frac{bh}{3}-\frac{ah}{3}+\frac{bc}{2}-\frac{ac}{2}\right]=\rho gh\left[\frac{ac}{2}+\frac{ah}{6}+\frac{bh}{3}+\frac{bc}{2}\right]. \]
答案：\(P=\rho gh\left(\frac{ac}{2}+\frac{ah}{6}+\frac{bh}{3}+\frac{bc}{2}\right)\)

8. 球从水中取出所做的功

题目：半径为 \(r\) 的球沉入水中，且与水面相切，其密度与水相同，计算把球从水中取出所做的功。

解答：球密度与水相同，所以浮力与重力平衡，只需考虑将球提出水面所做的功。

建立坐标系：以球心为原点，垂直向上为 \(z\) 轴正方向，水面在 \(z=r\) 处。球心初始位置 \(z=-r\)（球顶与水面相切），最终位置 \(z=r\)（球底与水面相切）。

在球心位置 \(z\) 时（\(z\in[-r,r]\)），露出水面的部分是一个球冠，高度为 \(h=r+z\)。记露出水面的部分体积为 \(V_{0}\)，球冠体积公式给出
\[ V_{0}=\pi h^{2}\left(r-\frac{h}{3}\right). \]
当球心在位置 \(z\) 时，露出水面的高度 \(h=r+z\)，所以
\[ V_{0}=\pi(r+z)^{2}\left(r-\frac{r+z}{3}\right)=\pi(r+z)^{2}\left(\frac{2r-z}{3}\right). \]
浮力减少量等于露出水面部分的水的重力：\(F=\rho gV_{0}\)。做功微元：\(dW=F\,dz=\rho gV_{0}dz\)。总功
\[ W=\int_{z=-r}^{r}\rho g\cdot\pi(r+z)^{2}\left(\frac{2r-z}{3}\right)dz. \]
令 \(u=r+z\)，则 \(z=u-r\)，\(dz=du\)，当 \(z=-r\) 时 \(u=0\)，当 \(z=r\) 时 \(u=2r\)。
\[ W=\frac{\pi\rho g}{3}\int_{0}^{2r}u^{2}\left(2r-(u-r)\right)du=\frac{\pi\rho g}{3}\int_{0}^{2r}u^{2}(3r-u)du \]
\[ =\frac{\pi\rho g}{3}\int_{0}^{2r}(3ru^{2}-u^{3})du=\frac{\pi\rho g}{3}\left[ru^{3}-\frac{u^{4}}{4}\right]_{0}^{2r}=\frac{\pi\rho g}{3}\left(r\cdot8r^{3}-\frac{16r^{4}}{4}\right)=\frac{\pi\rho g}{3}\left(8r^{4}-4r^{4}\right)=\frac{4}{3}\pi\rho gr^{4}. \]
答案：\(\frac{4}{3}\pi\rho gr^{4}\)
